\section{Билет 7. НЭП (новая экономическая политика) (1921--1928) в советской России}

\subsection{Причины}

Причины экономические:

\begin{enumerate}
    \item Упадок, вызванный Первой мировой и гражданскими войнами, интервенцией и политикой военного коммунизма.
    \item Продразвёрстка (обязательная сдача государству продукции сельского хозяйства).
    \item Уязвимое положение крестьянства в случае неурожая.
    \item Сокращение посевных площадей, поголовья скота.
    \item Национализация промышленности.
    \item Потеря доверия к деньгам.
    \item Возврат к натуральному хозяйству в деревне (выращивание товара только для собственных потребностей, не для продажи).
    \item Голод 1921--1923 и человеческие потери.
\end{enumerate}

\vspace{0.5cm}

Причины политические:

\begin{enumerate}
    \item Сопротивление политики военного коммунизма.
    \item Кронштадское востание (вся власть советам, а не .....).
\end{enumerate}

\vspace{0.5cm}

С 8 по 10 марта происходит X съезд РКП(б) (Российская коммунистическая партия большевиков) в дни Крондштатского восстания. Там закладываются меры НЭП (основы НЭПа):

\begin{enumerate}
    \item Замена продразвёрстки продналогом (система из 13 налогов в натуральной форме, чтобы заработать деньги на уплату налогов и приобретение продукций, крестьяне были вынуждены торговать с гос"=вом).
    \item Свободная торговля.
\end{enumerate}

\subsection{Основные меры. Сельское хозяйство:}

\begin{itemize}
    \item вводится продналог, который составлял определенный процент от выработки, который объявляется заранее, до начала посевного сезона;
    \item возврат возможности ведения единоличного хозяйства;
    \item в ограниченных размерах возвращены аренда земли и наёмный труд.
\end{itemize}

\subsection{Основные меры. Промышленность:}

\begin{itemize}
    \item разрешено частным лицам создавать мелкие и средние предприятия (кол"=во рабочих не больше 20);
    \item основной формой управления производством стали тресты --- объединения однородных или взаимосвязанных предприятий. Работая на условиях хоз расчета они самостоятельно решали что производить, где реализовывать продукцию и т.д. Законом предусматривалось, что "гос. казна за долги трестов не отвечает";
    \item возврат концессий --- разрешение правительства брать иностранным предприятиям предприятия в аренду;
    \item также вернулись совместные предприятия;
    \item поощерение коопераций (кооперативов);
    \item для предприятий появилась возможность хозрасчёта;
    \item гос. предприятия возвращаются к добровольному найму вместо трудовой повинности;
    \item стали приглашать старых и иностранных работников.
\end{itemize}

\subsection{Основные меры. Финансовая сфера:}

\begin{itemize}
    \item восстановление государственного банка и системы кредитных учреждений;
    \item в 1922--1924 гг. проводится денежная реформа, восстановившая денежную систему в стране к системе до Первой мировой войны, т.е. восстанавливается золотой стандарт. Вместо совзнаков вводятся медные и серебряные монеты, а также бумажные знаки;
    \item восстановлена система налогообложения, подоходный налог;
    \item вводится налог для нэпманов (частных предпринимателей), вводятся косвенные налоги;
    \item вводится плата за услуги: транспорта, связи, коммунального хозяйства.
\end{itemize}

\subsection{Основные меры. Торговля:}

\begin{itemize}
    \item востановлена частная торговля;
    \item востановлены крупные торговые сосредоточения: ярмарки, рынки;
    \item восстанавливаются товарные биржи.
\end{itemize}

\vspace{0.5cm}

Создаётся госплан как часть ВСНХ (Высший совет народного хозяйства), который планирует экономическую деятельность страны на будущее. 

Ввиду нехватки золота у государства в 1926 году был фактически отменён золотой стандарт.

Из"=за кризиса городов, многие люди переселяются в деревни, что сдерживает прогресс.

В 1920 году был утверждён план государственной комиссии по электрификации России (план ГОЭЛРО). Он подразумевал строительство нескольких электростанций. Но этот план мешает выполнять отсутствие ресурсов и специалистов. Поэтому была создана лишь одна электростанция --- Волховская.

\subsection{Финал НЭПа:}

\begin{enumerate}
    \item Советскому государству удалось за короткие сроки выйти из кризиса, восстановить промышленность и сельское хозяйство;
    \item В конечном итоге НЭП не смог решить основные экономические проблемы и власть начинает свёртывать НЭП. Официально НЭП был свернут 11 октября 1931 года. Начинается переход к индустриализации и коллективизации.
\end{enumerate}

\subsection{Чуть подробнее про свертывание НЭПа (причины свертывания):}

\begin{itemize}
    \item противоречия между административными и рыночными методами управления экономикой;
    \item ограничение участия частного капитала в экономике не позволяли использовать его для решения задачи индустриализации страны;
    \item усиление социального расслоения, появление нэпманов (новой буржуазии, эксплуататорских элементов) вызывало недовольство части населения;
    \item доминирование в обществе политической установки на временный характер НЭПа (если я правильно помню, народ реально сильно бухтел на тему того <<когда уже конец>> НЭПа);
    \item победа во внутриполитической борьбе 1920"=х гг. противников НЭПа, считавших, что НЭП был необходим только как временная мера для выхода из кризиса и что рыночные отношения несовместимы с социалистической идеей и практикой.
\end{itemize}


